\section{Discrete-in-space uniform controllability: the ODE-controllability bridge}
\label{sec:discrete}

This section covers the theory of \emph{space} semi-discretization -- the
technique actually used in this work -- and its relation to the
finite-dimensional (ODE) controllability theory of
Section~\ref{ssec:ode-prototype}. Internal (distributed) control is treated
as the primary case; boundary control appears only as a one-line contrast.

\subsection{Framing}
\label{ssec:discrete-framing}

\begin{quote}
  \itshape ``A space semi-discrete approximation of the parabolic control
  problems [\eqref{eq:heat-boundary} or \eqref{eq:heat-internal}] reduces the
  initial problem to a control problem for an ordinary differential equations
  (ODE's) system.''
\end{quote}
This reduction is the central bridge exploited throughout this section: study
of the semi-discrete heat equation's controllability becomes an application of
Theorems~\ref{thm:kalman}--\ref{thm:ode-hum} to a specific, mesh-dependent
family of linear ODE systems.

The numerical-HUM approach itself originates with the pioneering work of
\citet{glowinski1990numerical} for the wave equation, treated as an
unconstrained optimization problem for the adjoint system via convex-duality
theory \citep{ekeland1999convex}; it was later adapted to the heat equation to
find null controls by \citet{carthel1994exact}. HUM belongs to a broader
family of dual methods, including the \emph{weighted} HUM, which incorporates
a small perturbation to damp oscillatory control behavior
\citep{zuazua2006control}; more recently, primal methods working directly
with the primal system via Carleman weights \citep{fernandez2014numerical}
and least-squares/variational methods \citep{munch2014numerical} have given
results broadly comparable to the dual approach.

Discrete controllability of parabolic problems is usually studied via three
formulations -- time-discrete, space-discrete, and fully discrete
\citep{boyer2020carleman}. Uniform controllability (defined precisely in
Definition~\ref{def:uniform-controllability} below) is established
\emph{only} for the space-discrete formulation of the one-dimensional heat
equation \citep{zuazua2006control,lopez1998some}; the other two formulations
are not pursued in this work (see Section~\ref{ssec:other-formulations} for
why).

\subsection{The space-discrete ODE system}
\label{ssec:space-discrete-system}

Let $N \in \mathbb{N}$, let $h = L/(N+1)$ be the corresponding uniform mesh
size on $(0,L)$, $x_j = jh$ for $j=0,\dots,N+1$, and let
$y_h(t) = (y_1(t),\dots,y_N(t)) \in \R^N$ approximate
$\big(y(x_1,t),\dots,y(x_N,t)\big)$ via the standard central-difference
approximation of $\partial_{xx}$. For internal control, \eqref{eq:heat-internal}
becomes the linear ODE system on $\R^N$
\begin{equation}
  \label{eq:discrete-internal}
  y_h'(t) = A_h y_h(t) + B_h v_h(t), \qquad
  A_h = \frac{\alpha}{h^2}
  \begin{pmatrix}
    -2 & 1 & & \\
    1 & -2 & 1 & \\
    & \ddots & \ddots & \ddots \\
    & & 1 & -2
  \end{pmatrix} \in \R^{N\times N},
\end{equation}
where $v_h(t) \in \R^N$ and $B_h = \mathbf{1}_{\omega_h} \in \R^{N\times N}$ is
the diagonal matrix with $(\mathbf{1}_{\omega_h})_{jj}=1$ if $x_j \in \omega$
and $0$ otherwise (the identity matrix if $\omega$ is the whole domain). For
boundary control, the analogous system \citep[cf.][]{leveque2007finite} has
the same $A_h \in \R^{N\times N}$ but $v_h(t) \in \R$ a scalar and
$B_h \in \R^{N\times 1}$ a column vector, rather than $v_h(t) \in \R^N$ and
$B_h \in \R^{N\times N}$ as in \eqref{eq:discrete-internal}.

\begin{proposition}[Spectrum of $A_h$]
  \label{prop:spectrum}
  $A_h$ has eigenvalues
  $\lambda_j = \dfrac{2\alpha}{h^2}\big(1-\cos(j\pi h)\big)$, $j=1,\dots,N$,
  with eigenvectors $e^j_i = \sin(ij\pi h)$, $i=1,\dots,N$
  \citep[\S3.4]{leveque2007finite}.
\end{proposition}

This explicit, closed-form spectrum is precisely what makes the uniform
controllability results below tractable: the proof of
Theorem~\ref{thm:discrete-null-control} "relies solely on the fact that the
spectrum of the Laplacian can be computed explicitly."

\subsection{Uniform controllability}
\label{ssec:uniform-controllability}

\begin{definition}[Uniform controllability]
  \label{def:uniform-controllability}
  Controllability properties must be established via an observability
  inequality for the discrete adjoint system, rather than by directly
  invoking the Kalman rank condition, since the latter is a
  \textbf{qualitative, binary} criterion for a fixed finite-dimensional
  system: it says nothing about how the observability constant behaves as the
  system size grows with $h \to 0$. System \eqref{eq:discrete-internal} (resp.\
  its boundary analogue) is said to be \textbf{uniformly controllable with
  respect to $h$} if the discrete observability inequality below holds for a
  suitable constant $C>0$ \emph{independent of $h$}.
\end{definition}

The discrete adjoint system, with $\varphi_h(t) \in \R^N$, is
\begin{equation}
  \label{eq:discrete-adjoint}
  \varphi_h'(t) = -A_h^\top\varphi_h(t) \ \text{in } (0,T), \qquad
  \varphi_h(T) = \varphi_h^T \in \R^N.
\end{equation}

\begin{theorem}[Discrete observability inequality]
  \label{thm:discrete-observability}
  For every $T>0$ there is $C(T)>0$ such that
  \[
    h\sum_{j=1}^N |\varphi_j(0)|^2 \le C \int_0^T \left|\frac{\varphi_N(t)}{h}\right|^2 dt
  \]
  for every solution $\varphi_h$ of \eqref{eq:discrete-adjoint} and every
  $h>0$.
\end{theorem}
\begin{proof}[Source]
  \citet{lopez1998some}, based on a discrete observability inequality that
  relies solely on the explicit spectrum of Proposition~\ref{prop:spectrum}.
\end{proof}

\begin{theorem}[Discrete null controllability and convergence of controls]
  \label{thm:discrete-null-control}
  For any $T>0$ and $y_h^0$, there is a control $v_h \in L^2(0,T)$ steering
  the solution of \eqref{eq:discrete-internal} (or its boundary analogue) to
  $y_j(T)=0$, $j=1,\dots,N$. Moreover, for $y^0\in L^2(0,L)$, the controls
  $v_h$ may be built so that
  \begin{equation}
    \label{eq:convergence-controls}
    v_h \to v \ \text{in } L^2(0,T) \ \text{as } h \to 0,
  \end{equation}
  where $v$ is a null control for the continuous heat equation, provided the
  initial data are chosen appropriately.
\end{theorem}
\begin{proof}[Source]
  \citet{lopez1998some}.
\end{proof}

\begin{remark}[Extension to internal control]
  \label{rem:internal-extension}
  Theorem~\ref{thm:discrete-null-control} -- stated by \citeauthor{lopez1998some}
  for the boundary case -- can be modified to also give the convergence of the
  discretization for internal control, system \eqref{eq:discrete-internal}
  \citep[Rem.~3.4]{zuazua2006control}.
\end{remark}

\begin{remark}[Scope: only in 1D]
  \label{rem:1d-only}
  The uniform property of Definition~\ref{def:uniform-controllability} holds
  \emph{only for the one-dimensional} heat equation \citep{lopez1998some};
  convergence \eqref{eq:convergence-controls} follows from an explicit choice
  of $y^0$ via Fourier series \citep[Rem.~1.2]{lopez1998some}; and $v_h$
  converges to the minimal-norm HUM control of
  Section~\ref{ssec:ode-prototype}. In two dimensions, the space semi-discrete
  heat equation is \textbf{not} controllable at all \citep{zuazua2005propagation}
  -- a counterexample that is dimension-general, not specific to internal or
  boundary control, and that directly explains the 1D restriction here.
  Further references establishing/refining uniform controllability for the
  space-discrete 1D case, not worked through in detail here, include
  \citet{labbe2006uniform} and \citet{boyer2010discrete}.
\end{remark}

\begin{remark}[Gap, not filled: joint $(h,k)$ scaling]
  \label{rem:hk-gap}
  Every uniform-controllability result above is for \emph{space}
  semi-discretization only ($h\to0$, time kept continuous); no joint
  $(h,k)$-refinement (fully discrete) uniform-controllability theorem is
  established here. This is the discrete-side counterpart of the open
  condition-number scaling question already noted after
  Proposition~\ref{prop:condition-number}: the penalized-HUM condition number
  $\nu_a = 1+\|\Lambda\|/\varepsilon$ is stated purely as a function of
  $\varepsilon$, with no established $(h,k)$-refinement scaling law. Both gaps
  should be read together, and neither should be presented as resolved.
\end{remark}

\subsection{Other formulations (brief, not pursued here)}
\label{ssec:other-formulations}

The time-discrete and fully-discrete formulations are noted only for
completeness, since this work uses space discretization exclusively.
\textbf{Time-discrete}: this is a genuine \emph{negative} result, not a milder
version of the space-discrete case above -- the time-semi-discrete
multi-dimensional heat equation is in general neither null nor approximately
controllable \citep{zheng2008controllability}, and Carleman estimates for the
time-discrete operator give only "another sense of controllability"
\citep{boyer2020carleman}, not the same uniform-observability property.
\textbf{Fully discrete}: results are scarce; null controllability for a
family of one-dimensional parabolic equations has been shown via Carleman
inequalities, again only in a relaxed sense of controllability
\citep{casanova2021carleman}.
